%%%%

%%%   Самарский государственный технический университет

%%%   Кафедра прикладной математики и информатики

%%%

%%%   Преамбула для статьи в журнал "Вестник Самарского государственного

%%%   технического университета. Серия: «Физико-математические науки»"

%%%   Ориентирована под MikTEx 2.4 for Win32

%%%

%%%   Хотя сам журнал собирается под GNU/Linux на дистрибутиве TeXLive



\documentclass[11pt,a4paper,twoside]{article}



\usepackage[cp1251]{inputenc}

\usepackage[T2A]{fontenc}

\usepackage[english,russian]{babel}

\usepackage{samgtubib}

\usepackage[dvips]{graphicx}

\usepackage[multiple]{footmisc}



\usepackage{samgtu}

\renewcommand{\VestnikYear}{2009} % Год издания Вестника

\renewcommand{\VestnikNumber}{2\,(19)} % Номер Вестника

\renewcommand{\VestnikMonth}{Ноябрь} % Месяц выхода Вестника

\renewcommand{\VestnikData}{01.11.09} % Дата подписания Вестника в печать

\renewcommand{\VestnikUPL}{26,64} % Усл. печ. л.

\renewcommand{\VestnikUIL}{25,73} % Уч.-изд. л.

\renewcommand{\VestnikRegNumber}{479/09} % Регистрационный номер

\renewcommand{\VestnikOrder}{994} % Номер заказа







%

%

%\usepackage{qrcode}

%%

%%\newcommand{\totop}[1]{\raisebox{6pt}[1pt][2pt]{#1}} 

%%\newcommand{\tottop}[1]{\raisebox{12pt}[1pt][2pt]{#1}} 

%%\newcommand{\totttop}[1]{\raisebox{18pt}[1pt][2pt]{#1}} 

%%\newcommand{\tottttop}[1]{\raisebox{24pt}[1pt][2pt]{#1}} 

%%\newcommand{\totttttop}[1]{\raisebox{30pt}[1pt][2pt]{#1}} 

%\newcommand{\tobot}[1]{\raisebox{-6pt}[1pt][2pt]{#1}} 

%%\newcommand{\tobbot}[1]{\raisebox{-12pt}[1pt][2pt]{#1}} 

%%\newcommand{\tobbbot}[1]{\raisebox{-18pt}[1pt][2pt]{#1}}

%

%\graphicspath{{./00PIC/}}

%

%%\samgtupubl % для генерации pdf, отдаваемого в типографию СамГТУ

%\pdfcrop % обрезка pdf для размещения на math-net.ru

%%\draft % На корректуре

%\filllastpage % Последняя страница не заполнена не полностью

%%\briefcomm % Статья является кратким сообщением



\vsgtu{1452}



\FirstPage{7}

\LastPage{21}

%

%

%\begin{document}





\renewcommand{\baselinestretch}{0.9}



%\hfill \qrcode[hyperlink,height=.8in]{http://dx.doi.org/10.14498/vsgtu1452}

%\vspace{-.8in}







%%%%%%%%%%%%%%%%%%%%%%%%%%%%%%%%%%%%%%%%%%%%%%%%%%%%%%%%%%%%%%%%%%%%%%%%%%%%%%%%%%%%

%Информация о статье и об авторах на русском и английском языках



\udc{517.956.3}

\msc{35L80}% Коды MSC см. http://www.ams.org/msc/



\rutitle{Задача Гурса для нагруженного вырождающегося~гиперболического уравнения~второго порядка с~оператором Геллерстедта в~главной части}% Полный заголовок

{Задача Гурса для нагруженного вырождающегося гиперболического уравнения \dots}



\entitle{Goursat problem for loaded degenerate second order hyperbolic equation with Gellerstedt operator in principal part}



\ruauthor{А.~Х.~Аттаев}{А\,т\,т\,а\,е\,в~А.~Х.}

\enauthor{A.~H.~Attaev}{A\,t\,t\,a\,e\,v~A.~H.}



\ruaffil{\runiipma.}

\enaffil{\enniipma.}





\ruauthordetails{1}{\fullauthor{\rupid{45430}{Анатолий Хусеевич Аттаев}}\regalia{к.ф.-м.н., доц.; \mem{attaev.anatoly@yandex.ru}}\position{заведующий отделом} \dept{отдел САПР смешанных систем и управления} }





\enauthordetails{1}{\fullauthor{\enpid{45430}{Anatoly H. Attaev}} \regalia{Cand. Phys. \& Math. Sci.; \mem{attaev.anatoly@yandex.ru}} \position{Head of Dept.}

\dept{Dept. CAD of Mixed Systems and Management} }







\enkeywords{anomalous diffusion, diffusion fractional order, Riemann--Liouville operator, fundamental solution, general representation of solution, modified Bessel function, Wright function, integral transformation wich Wright function in kernel}









\rukeywords{задача Гурса, нагруженное уравнение, гиперболическое уравнение, вырождающееся уравнение, 

оператор Геллерстедта, метод функции Грина--Адамара}



\ruabstract{Рассматривается нагруженное вырождающееся гиперболическое уравнение второго порядка с~переменными коэффициентами. Главная часть уравнения представляет собой оператор Геллерстедта. 

Нагруженное слагаемое представляет собой след искомого решения на линии вырождения, которая лежит внутри области. 

Исследуется задача с~данными на одной из характеристик исследуемого уравнения. В~модельном случае, когда коэффициенты при младших членах обращаются в ноль, решение задачи Гурса выписано в~явном виде. 

При этом использовалась функция Грина"--~Адамара для уравнения Эйлера"--~Дарбу"--~Пуассона. 

В~общем случае разрешимость задачи Гурса эквивалентным образом редуцирована к вопросу о разрешимости интегрального уравнения Вольтерра второго рода.  

В~этом случае использована схема, реализованная С.~Геллерстедтом при доказательстве существования решения второй задачи Дарбу для рассматриваемого уравнения без нагрузки. 

В~обоих случаях существенно использовались известные свойства функции Грина"--~Адамара.}







\enabstract{In the paper we study a loaded degenerate hyperbolic equation of the second order with variable coefficients. The principal part of the equation is the Gellerstedt operator. The loaded term is given in the form of the trace of desired solution on the degenerate line. The latter is located in the inner part of the domain. We investigate a boundary value problem. The boundary conditions are given on a characteristics line of the equation under study. For the model equation (when all subordinated coefficients are zero) we construct an explicit representation for solution of the Goursat problem. In the general case, we reduce the problem to an integral Volterra equation of the second kind. We apply the method realized by Sven Gellerstedt solving the second Darboux problem. In both cases, model and general, we use widely properties of the Green--Hadamard function.}



\enkeywords{Goursat problem, loaded equation, hyperbolic equation, degenerate equation, Gellerstedt operator, the Green--Hadamard's function method}







\date{13/X/2015} % Дата поступления статьи в редакцию.

\revisiondate{23/X/2015} % В окончательном виде

\accepteddate{24/XI/2015}







\makerutitle



%%%%%%%%%%%%%%%%%%%%%%%%%%%%%%%%%%%%%



\newpage





Исследованию краевых задач для линейных нагруженных дифференциальных уравнений второго порядка с двумя независимыми переменными посвящено много работ. 

Достаточно полная библиография приведена в монографии \cite{litAtt01}.

Нагруженным уравнением в области $\Omega\in\mathbb R^n$ называется уравнение 

$$

L u(x)=f(x),

$$ 

где $L$ "---  оператор, содержащий след некоторых операций от искомого решения $u(x)$ на принадлежащих $\bar{\Omega}$ многообразиях размерности меньше $n$~\cite{litAtt02}.



Существенную роль для корректной разрешимости тех или иных задач для нагруженных уравнений играет не просто наличие следа $\Gamma$ от искомого решения, а в большей степени следообразующее отображение 

$$F:\bar{\Omega}\rightarrow \Gamma.$$ 

Особенно этот эффект проявляется для нагруженных гиперболических уравнений. 

Поясним это на примере. 

Пусть заданы два модельных уравнения 

$$

u_x+u_y=\lambda u(x, 0)

\qquad\textup{и}\qquad

u_x+u_y=\lambda u(x-y, 0),

$$

где $\lambda=\rm const$.

След от искомого решения в правых частях есть одно и то же многообразие, хотя легко заменить, что общие интегралы этих двух уравнений принципиально отличаются.



Задачам с данными на характеристических многообразиях для нагруженных строго и слабо гиперболических уравнений посвящены работы~\mbox{[\citen{litAtt03}--\citen{litAtt07}].} 

Характеристической задаче и~задаче Дарбу для широкого класса систем нагруженных вырождающихся гиперболических уравнений посвящены работы \cite{litAtt08, litAtt09}.



В данной работе объектом исследования является уравнение вида

\begin{multline}

\label{Attfor2015_1}

L^m u=u_{yy}-|y|^m u_{xx}+a(x,y)u_x+b(x,y)u_y+c(x,y)u=

\\

=\lambda\sum\limits_{i=1}^n a_i (x,y) D_{0\xi}^{\alpha_i} u(\xi , 0)+d(x,y),

\quad \xi=x-\frac{2}{m+2}|y|^{\frac{m+2}{2}}.

\end{multline}



Пусть $\Omega$ "--- область, ограниченная характеристиками 

\begin{equation}

\label{Attfor2015_2}

x-\frac{2}{m+2}|y|^{\frac{m+2}{2}}=0, \quad x+\frac{2}{m+2}|y|^{\frac{m+2}{2}}=1

\end{equation}

уравнения \eqref{Attfor2015_1}. 

Здесь $\lambda$ и $m$ "--- действительные числа, 

причем $$m\geqslant 0,\quad \alpha_1<\alpha_2<\ldots<\alpha_n=\alpha<1;$$ 

$D_{a x}^l$  "--- оператор дробного интегрирования (при $l<0$) и дробного дифференцирования 

(при $l>0$) порядка $|l|$ с~началом в точке $a$, определяемый как и~в~\cite{Attfor2015_lit_01} формулой

\begin{equation}

D_{a x}^l f(t) = \left\{

\begin{array}{cl}

\displaystyle \frac{\mathop{\rm sign}(x-a)}{\Gamma(-l)}\int_a^x \frac{f(t)dt}{|x-t|^{l+1}}, 

& l<0,\\[5mm]

 f(x),                                                                   

& l=0,\\[3mm]

\displaystyle {\rm sign}^{[l]+1}(x-a)\frac{d^{[l]+1}}{dx^{[l]+1}} D_{a x}^{l-[l]-1} f(t),           

& l>0,

\end{array} \right.

\end{equation}

где $[l]$ "--- целая часть $l$, $[l]\leqslant l<[l]+1$, 

$\Gamma(z)$ "--- Гамма"=функция Эйлера.



\smallskip



{\small \sc Задача Гурса.} 

\it 

Найти регулярное в области $\Omega$ при $y\neq 0$ 

решение $u(x,y)$ уравнения \eqref{Attfor2015_1} из класса

$C(\bar{\Omega})\cap C^1(\Omega),$ удовлетворяющее условию

\begin{equation}

\label{Attfor2015_3}

u\left(\frac{x}{2}, \mathop{\rm sign} y \Bigl(\frac{m+2}{4} x\Bigr)^\frac{2}{m+2}\right)

=\varphi^{\pm} (x), \quad 0\leqslant x \leqslant 1,

\end{equation}

где знак ``$+$'' берется при $y>0$ и знак ``$-$'' "--- при $y<0$.









Предполагается$,$ что $\varphi^+(0)=\varphi^-(0);$

$

b,$ $ c,$ $ d,$  $a_i\in C^1(\bar{\Omega})\cap C^3(\Omega),$

$i=1, 2, \dots, n;$ 

\begin{equation}

\label{Attfor2015_4}

\varphi^\pm \in C^1(\bar{J})\cap C^3(J),

\end{equation}

где $ \bar J =\{( x, y ):  0 \leqslant x \leqslant 1, y=0\}.$







\smallskip



\rm



Имеет место следующая теорема.



\smallskip



\hypertarget{att:th1}{}



{\small\sc Теорема 1.} \it  

Пусть коэффициенты $a(x, y),$ $c(x, y)$ 

по переменной $y$ являются четными функциями$,$ 

а~коэффициент $b(x, y)$ "---

нечетной$;$ 

коэффициент

$a(x,y)$  удовлетворяет условию Геллерстедта\/{\rm:}

\begin{equation}

\label{Attfor2015_5}

a(x, y)=a_0(x,y)|y|^{\gamma (m)}, \quad  

a_0\in C^1(\bar{\Omega})\cap C^3(\Omega),

\end{equation}

где $\gamma(m)=0$ при $m<2$ и $\gamma(m)={\rm const }>{m}/{2}-1$ при $m\geqslant 2.$



Тогда задача Гурса разрешима и притом  единственным образом.



\rm



\smallskip



Сначала приведем \textit{д\,о\,к\,а\,з\,а\,т\,е\,л\,ь\,с\,т\,в\,о\/} теоремы \hyperlink{att:th1}{1}  для случая

$a=b=c=d\equiv 0$, $a_1=1$, $\alpha_1=0$; $a_i(x,y)=0$, $i=2, 3, \dots, n$,

т.е. для модельного локально нагруженного уравнения Геллерстедта:

\begin{equation}

\label{Attfor2015_6}

u_{yy}-|y|^m u_{xx}=\lambda u\Bigl( x-\frac{2}{m+2}|y|^{\frac{m+2}{2}} , 0 \Bigr).

\end{equation}







Через $\Omega^+$ и $\Omega^-$ обозначим части $\Omega$, расположенные в~верхней  и~нижней полуплоскостях соответственно.



В области $\Omega^-$ уравнение \eqref{Attfor2015_6} и~краевые условия \eqref{Attfor2015_3} имеют вид

\begin{gather}

\label{Attfor2015_7}

u_{yy}-(-y)^m u_{xx}=\lambda u\Bigl( x-\frac{2}{m+2}(-y)^{\frac{m+2}{2}} , 0 \Bigr),

\\

\label{Attfor2015_8}

u\Bigl(\frac{x}{2}, - \Bigl(\frac{m+2}{4} x\Bigr)^\frac{2}{m+2}\Bigr)

=\varphi^{-} (x), \quad 0\leqslant x \leqslant 1.

\end{gather}

Соотношения \eqref{Attfor2015_7} и \eqref{Attfor2015_8} в характеристических переменных

\begin{equation}

\label{Attfor2015_9}

\xi=x-\frac{2}{m+2}(-y)^{\frac{m+2}{2}}, \quad \eta=x+\frac{2}{m+2}(-y)^{\frac{m+2}{2}}

\end{equation}

записываются в виде

\begin{equation}

\label{Attfor2015_10}

v_{\xi\eta}+\frac{\beta}{\eta-\xi}(v_\eta-v_\xi)=\frac{\mu v(\xi,\xi)}{(\eta-\xi)^{4\beta}},

\end{equation}

\begin{equation}

\label{Attfor2015_11}

v(0,\eta)=\psi(\eta)\equiv \varphi^-(\eta),\quad 0\leqslant\eta\leqslant 1,

\end{equation}

где $4\mu=\lambda(2-4\beta)^{4\beta}$, $\beta={m}/({2m+4})$,

$$

v(\xi,\eta)=u\Bigl(\frac{\xi+\eta}{2},

-\Bigl(\frac{m+2}{4} \Bigr)^{1-2\beta}(\eta-\xi)^{1-2\beta}\Bigr).

$$



Пусть $H(\xi,\eta,\xi_0,\eta_0)$ "--- функция Грина"--~Адамара

 второй задачи Дарбу для уравнения Эйлера"--~Дарбу"--~Пуассона

\begin{equation}

\label{Attfor2015_12}

Lv\equiv v_{\xi\eta}-\frac{\beta}{\eta-\xi}(v_\eta-v_\xi)=0,

\end{equation}

определяемая следующими условиями:

\begin{itemize}

\item[1)] $H(\xi, \eta, \xi_0, \eta_0)$ как функция от  $(\xi,\eta)$

 есть решение сопряженного уравнения  $L^*v=0$, 

 а как функция от  $(\xi_0,\eta_0)$ "--- решение уравнения \eqref{Attfor2015_10};



\item[ 2)] $H(\xi_0,\eta_0,\xi_0,\eta_0)=1$;



\item[ 3)] $H(\xi,\xi,\xi_0,\eta_0)=0$;



\item[ 4)] $\displaystyle \frac{\partial[H]}{\partial\xi}+\beta\frac{[H]}{\eta-\xi}=0$ при $\eta=\xi_0$, где

$$

[H]=\lim\limits_{\varepsilon\rightarrow 0}\bigl(H(\xi,\xi_0+\varepsilon,\xi_0,\eta_0)

-H(\xi,\xi_0-\varepsilon,\xi_0,\eta_0)\bigr).

$$



\end{itemize}





Как известно, функция Грина"--~Адамара $H(\xi,\eta,\xi_0,\eta_0)$ 

была явно построена Геллерстедтом \cite{Attfor2015_lit_02} и имеет вид

\begin{displaymath}

H(\xi,\eta;\xi_0,\eta_0) = 

\left\{ \begin{array}{ll}

R(\xi,\eta;\xi_0,\eta_0), \quad \eta\geqslant\xi_0,\\

\bar{R}(\xi,\eta;\xi_0,\eta_0), \quad \eta\leqslant\xi_0,

\end{array} \right.

\end{displaymath}

где $R(\xi,\eta;\xi_0,\eta_0)=(\eta-\xi)^\beta(\eta_0-\xi_0)^{-\beta}

F(\beta , 1-\beta, 1, \sigma)$ "--- функция Римана,

$$

\bar{R}(\xi,\eta;\xi_0,\eta_0)=

\gamma\frac{(\eta-\xi)^{2\beta}}{(\xi_0-\xi)^\beta(\eta_0-\eta)^\beta}

F\left(\beta,\beta, 2\beta; {1}/{\sigma}\right).

$$

Здесь

$$

\gamma=\frac{\Gamma(\beta)}{\Gamma(1-\beta)\Gamma(2\beta)},\quad 

\sigma=\frac{(\xi_0-\xi)(\eta_0-\eta)}{(\eta-\xi)(\eta_0-\xi_0)}.

$$



Пусть существует решение задачи Гурса \eqref{Attfor2015_3} для 

уравнения \eqref{Attfor2015_6} и пусть функция

$$

\nu(x)=u_y(x , 0),\quad 0<x<1

$$

имеет дробный интеграл порядка $1-2\beta$ с~началом в точке $0$ и~с~концом в~любой точке $x\in{}]0 , 1]$. 

Тогда, пользуясь свойством функции $H$, по схеме, реализованной Геллерстедтом \cite{Attfor2015_lit_02},

при доказательстве существования решения второй задачи Дарбу можно доказать, что 

\begin{multline}

\label{Attfor2015_13}

u(x,y)=\frac{1}{2}\Bigl(\frac{4}{m+2} \Bigr)^{2\beta}\gamma

 \int _0^\xi \nu(t)(\xi-t)^{-\beta}(\eta-t)^{-\beta}dt+

\\

+ \int _0^\eta \Bigl( \psi'(t)+\frac{\beta}{t}\psi(t) \Bigr)

H(0, t;\xi,\eta) dt + \\

+\mu \int _0^\xi d\xi_1  \int _{\xi_1}^\eta

\frac{u(\xi_1 , 0)}{(\eta_1-\xi_1)^{4\beta}} H(\xi_1,\eta_1;\xi,\eta) d\eta_1.

\end{multline}







Из \eqref{Attfor2015_13} при $y=0$, $\eta=\xi=x$, $t<\xi$, 

с учетом пределов

$$

\lim\limits_{\eta-\xi\rightarrow +0} H(0,t;\xi,\eta)=

\lim\limits_{\eta-\xi\rightarrow +0} \bar{R}(0,t;\xi,\eta)=

\gamma t^{2\beta}\xi^{-\beta} (\xi-t)^{-\beta}, \quad t<\xi;

$$

\begin{multline}

\lim\limits_{\eta-\xi\rightarrow +0} H(\xi_1,\eta_1;\xi,\eta)=

\lim\limits_{\eta-\xi\rightarrow +0} \bar{R}(\xi_1,\eta_1;\xi,\eta)= \\ =

\gamma (\eta_1-\xi_1)^{2\beta} (\xi-\eta_1)^{-\beta} (\xi-\xi_1)^{-\beta}, \quad \eta_1<\xi,

\end{multline}

имеем

\begin{multline}

\label{Attfor2015_14}

u(x , 0)=\frac{1}{2}\Bigl(\frac{4}{m+2} \Bigr)^{2\beta}\gamma

 \int _0^x \frac{\nu(t)dt}{(x-t)^{2\beta}}+

\\

+\mu\gamma \int _0^x d\xi_1  \int _{\xi_1}^x

\frac{u(\xi_1 , 0)\,d\eta_1}{(\eta_1-\xi_1)^{2\beta}(x-\eta_1)^\beta(x-\xi_1)^\beta} +

\\

+\gamma x^{-\beta} \int _0^x \bigl(t\psi'(t)+\beta\psi(t)\bigr) t^{2\beta-1}(x-t)^{-\beta}dt.

\end{multline}



Произведем замену переменной интегрирования $\eta_1$ по формуле 

$$\eta_1=\xi_1+(x-\xi_1)t.$$

Тогда, принимая во внимание, что

\begin{multline}

 \int _{\xi_1}^x (x-\eta_1)^{-\beta}(\eta_1-\xi_1)^{-2\beta}d\eta_1= \\ =

(x-\xi_1)^{1-3\beta} \int _0^1 t^{-2\beta}(1-t)^{-\beta}dt=

(x-\xi_1)^{1-3\beta}B(1-2\beta , 1-\beta)

\end{multline}

и используя определение дробного интеграла, соотношению \eqref{Attfor2015_14} 

можно придать более компактный вид:

\begin{equation}

\label{Attfor2015_15}

\tau(x)=\lambda\gamma_1 D_{0x}^{4\beta-2}\tau(t)+

\gamma_2 D_{0x}^{2\beta-1}\nu(t)+\Psi(x), \quad 0<x<1

\end{equation}

где $$\tau(x)\equiv u(x , 0), \qquad

\gamma_1=\frac{1}{4}(2-4\beta)^{4\beta}

\gamma B(1-2\beta , 1-\beta)\Gamma(4\beta),

$$

$$

\gamma_2=\frac{1}{2}\Bigl(\frac{4}{m+2}\Bigr)^{2\beta}

\gamma \Gamma(2-4\beta),

$$

$$

\Psi(x)=\gamma\Gamma(1-\beta)x^{-\beta} D_{0x}^{\beta-1}

\bigl( t\psi'(t)+\beta\psi(t) \bigr)t^{2\beta-1}.

$$



Равенство \eqref{Attfor2015_15} представляет собой фундаментальное 

соотношение между функциями $\tau(x)$ и $\nu(x)$, принесенное 

из области $\Omega^-$ на линию $y=0$ параболического вырождения.

Найдем теперь соотношение, которое приносится из области $\Omega^+$.



В области $\Omega^+$  имеем

\begin{gather}

\label{Attfor2015_16}

u_{yy}-y^m u_{xx}=\lambda  u\Bigl( x-\frac{2}{m+2} y^{\frac{m+2}{2}} , 0\Bigr) ,\quad y>0,

\\

\label{Attfor2015_18}

u\Bigl(\frac{x}{2},\Bigl( \frac{m+2}{4} x\Bigr)^{\frac{2}{m+2}} \Bigr)=\varphi^+(x) ,\quad 0\leqslant x\leqslant 1,

\\[2mm]

\label{Attfor2015_18}

\tau(x)= u(x , 0),\quad  \lim\limits_{y\rightarrow 0+}u_y(x,y)=\nu(x) ,\quad 0<x<1.

\end{gather}



В системе \eqref{Attfor2015_16}--\eqref{Attfor2015_18} произведем преобразование 

$$\widetilde{x}=x,\quad \widetilde{y}=-y,\quad \widetilde{u}=\widetilde{u}(\widetilde{x},\quad \widetilde{y})=u(x,-y).$$

После этого она примет вид

\begin{gather}

\widetilde{u}_{\widetilde{y}\widetilde{y}}-(-\widetilde{y})^m 

\widetilde{u}_{\widetilde{x}\widetilde{x}}=

\lambda  \widetilde{u}\Bigl( \widetilde{x}-\frac{2}{m+2} 

(-\widetilde{y})^{\frac{m+2}{2}} , 0\Bigr) ,\quad \widetilde{y}<0,

\\

\widetilde{u}\Bigl(\frac{\widetilde{x}}{2},-\Bigl( \frac{m+2}{2}

 \widetilde{x}\Bigr)^{\frac{2}{m+2}} \Bigr)=\varphi^+(\widetilde{x}) 

 ,\quad 0\leqslant \widetilde{x}\leqslant 1,

\\[2mm]

\widetilde{u}(\widetilde{x} , 0)=\tau(\widetilde{x}) , \quad 

 \lim\limits_{\widetilde{y}\rightarrow 0-}\widetilde{u}_{\widetilde{y}}=-\nu(x) ,\quad 0<x<1.

\end{gather}

Поэтому нетрудно заключить, что фундаментальное соотношение

между $\tau(\widetilde{x})$ и $\nu(\widetilde{x})$,

 принесенное из области $\Omega^+$, задается формулой

\begin{equation}

\label{Attfor2015_19}

\tau(x)=\lambda\gamma_1 D_{0x}^{4\beta-2}\tau(t)-

\gamma_2 D_{0x}^{2\beta-1}\nu(t)+\widetilde{\Psi}(x), \quad 0<x<1,

\end{equation}

где

$$

\widetilde{\Psi}(x)=\gamma\Gamma(1-\beta)x^{-\beta} D_{0x}^{\beta-1}

\bigl(t\widetilde{\psi}'(t)+\beta\widetilde{\psi}(t)\bigr)t^{2\beta-1},

$$

$$

\widetilde{\psi}(t)=\varphi^+(t), \quad 0\leqslant t \leqslant 1.

$$



Cистема \eqref{Attfor2015_15}, \eqref{Attfor2015_19}, очевидно, эквивалентна системе

\begin{gather}

\label{Attfor2015_20}

D_{0x}^{2\beta-1}\nu(t)=f_1(x),\quad 0<x<1,

\\

\label{Attfor2015_21}

\tau(x)-\lambda\gamma_1 D_{0x}^{4\beta-2}\tau(t)=f_2(x),\quad  0\leqslant x \leqslant 1,

\end{gather}

где

\begin{gather}

\label{Attfor2015_22}

f_1(x)=\frac{1}{2\gamma_2}\left[\Psi(x)-\widetilde{\Psi}(x)\right],

\\

\label{Attfor2015_23}

f_2(x)=\frac{1}{2}\left[\Psi(x)+\widetilde{\Psi}(x)\right].

\end{gather}



Итак, доказана лемма, которая очень важна для доказательства теоремы~\hyperlink{att:th1}{1}

в~случае уравнения \eqref{Attfor2015_6}.



\smallskip





{\small \sc Лемма.} \it  

Решение $u(x,t)$ задачи Гурса \eqref{Attfor2015_3} для уравнения \eqref{Attfor2015_6} в области $\Omega$  существует тогда и только тогда$,$ когда функция $\nu(x)=u_y(x , 0)$

является решением уравнения Абеля \eqref{Attfor2015_20}

с правой частью \rm \eqref{Attfor2015_22}, \it  а $\tau(x)=u(x , 0)$ "---

решение уравнения \eqref{Attfor2015_21} с~правой частью  \rm \eqref{Attfor2015_23}.





\rm 

\smallskip





Исследуем разрешимость уравнений \eqref{Attfor2015_20} и  \eqref{Attfor2015_21}. 

С~этой целью изучим свойства функции  $f_1(x)$ и $f_2(x)$.



Из \eqref{Attfor2015_4} заключаем, что 

\begin{equation}

\label{Attfor2015_24}

\Phi(x)=x\psi'(x)+\beta\psi(x) \in C(\bar{J})\cap C^2(J).

\end{equation}

Ясно, что

$$

\Psi(x)=\gamma\Gamma(1-\beta)x^{-\beta}D_{0x}^{\beta-1}

t^{2\beta-1}\Phi(t)=\gamma x^{-\beta} \int _0^x

\frac{\Phi(t)t^{2\beta-1}}{(x-t)^\beta} dt.

$$

Отсюда после замены $t=xz$ находим

\begin{equation}

\label{Attfor2015_25}

\Psi(x)=\gamma \int _0^1 z^{2\beta-1}(1-z)^{-\beta}\Phi(xz)dz.

\end{equation}

Стало быть, $\Psi(x)\,\in C(\bar{J})\cap C^2(J)$. 

Теперь нет сомнений в том, что 

\begin{equation}

\label{Attfor2015_26}

f_1(x), \, f_2(x)\in C(\bar{J})\cap C^2(J).

\end{equation}

В силу \eqref{Attfor2015_25} решение уравнения Абеля \eqref{Attfor2015_20} имеет вид

\begin{equation}

\label{Attfor2015_27}

\nu(x)=D_{0x}^{1-2\beta}f_1(t).

\end{equation}

Из \eqref{Attfor2015_25} c учётом \eqref{Attfor2015_24} видно, что

$$

\Psi(0)=\gamma\Phi(0) \int _0^1 z^{2\beta-1}(1-z)^{-\beta}dz=\gamma\beta\psi(0) B(2\beta , 1-\beta).

$$

Так как

$$

\gamma \beta B(2\beta , 1-\beta)=\frac{\Gamma(\beta)\beta}{\Gamma(2\beta)\Gamma(1-\beta)}

\frac{\Gamma(2\beta)\Gamma(1-\beta)}{\Gamma(1+\beta)}=1

$$

и $\psi(0)=\varphi(0)$, получаем $\Psi(0)=\varphi(0)$. 

Аналогично, $\widetilde{\Psi}(0)=\varphi(0)$.

Следовательно, $f_1(0)=0$. 

С~учетом этого \eqref{Attfor2015_27} можно переписать в виде

\begin{multline}

\nu(x)=\frac{d}{dx}D_{0x}^{-2\beta}f_1(x)=

\frac{d}{dx}\frac{1}{\Gamma(2\beta)} \int _0^x

f_1(t)(x-t)^{2\beta-1}dt=

\\

=\frac{d}{dx}\frac{1}{\Gamma(2\beta)}\biggl( 

-\frac{1}{2\beta} f_1(t)(x-t)^{2\beta}\Big|_0^x

+\frac{1}{2\beta} \int _0^x f_1'(t)(x-t)^{2\beta}dt

\biggr)=

\\

=\frac{1}{2\beta\Gamma(2\beta)}\frac{d}{dx} \int _0^x f_1'(t)(x-t)^{2\beta}dt=

\frac{1}{\Gamma(2\beta)} \int _0^x \frac{f_1'(t)dt}{(x-t)^{1-2\beta}},

\end{multline}

т.~е.

\begin{equation}

\label{Attfor2015_28}

\nu(x)=D_{0x}^{-2\beta}f_1'(t).

\end{equation}



Из \eqref{Attfor2015_28} и \eqref{Attfor2015_26} вытекает, что функция $\nu(x)\in C(\bar{J})\cap C^1(J)$

и при $x\rightarrow 0$ обращается в нуль.



Поскольку функция $f_2(x)\in C(\bar{J})\cap C^2(J)$, решение уравнения Вольтерра

второго рода \eqref{Attfor2015_21} принадлежит $C(\bar{J})\cap C^2(J)$.



Доказана (см. формулу \eqref{Attfor2015_13}) следующая теорема.



\smallskip



{\small \sc Теорема 2.} \it  Единственное решение задачи Гурса \eqref{Attfor2015_3}

для уравнения \eqref{Attfor2015_6} в~области $\Omega^-$ задается формулой

\begin{multline}

u(x,y)=\frac{1}{2} \Bigl(\frac{4}{m+2} \Bigr)^{2\beta} \gamma  \int _0^\xi

\frac{D_{0t}^{-2\beta}f_1'(t)dt}{\bigl((\xi-t)(\eta-t)\bigr)^\beta}+

 \int _0^\eta \frac{\Phi(t)}{t}H(0,t;\xi,\eta)dt+

\\

+\mu \int _0^\xi d\xi_1  \int _{\xi_1}^\eta

\frac{\tau(\xi_1)}{(\eta_1-\xi_1)^{4\beta}}

H(\xi_1,\eta_1;\xi,\eta)d\eta_1.

\end{multline}





\rm



\smallskip



Вернемся к случаю уравнения \eqref{Attfor2015_1}. 

В характеристических координатах \eqref{Attfor2015_9} оно приводится к виду

\begin{multline}

\label{Attfor2015_29}

v_{\xi\eta}+\frac{\beta}{\eta-\xi}(v_\xi-v_\eta)-(\eta-\xi)^{-4\beta}A(\xi,\eta)(v_\xi+v_\eta)-

\\

-(\eta-\xi)^{-2\beta}B(\xi,\eta)(v_\xi-v_\eta)-(\eta-\xi)^{-4\beta}C(\xi,\eta)v=

\\

=(\eta-\xi)^{-4\beta}\bigl(A_i(\xi,\eta)D_{0\xi}^{\alpha_i}v(\xi,\xi)+f(\xi,\eta)\bigr),

\end{multline}

где

$$

f(\xi,\eta)=-\frac{1}{4} \Bigl(\frac{4}{m+2}\Bigr)^{4\beta} d

\Bigl(\frac{\xi+\eta}{2},-\Bigl(\frac{m+2}{4}\Bigr)^{1-2\beta}

(\eta-\xi)^{1-2\beta} \Bigr),

$$

$$

A_i(\xi,\eta)=\frac{1}{4} \Bigl(\frac{4}{m+2}\Bigr)^{4\beta} a_i

\Bigl(\frac{\xi+\eta}{2},-\Bigl(\frac{m+2}{4}\Bigr)^{1-2\beta}

(\eta-\xi)^{1-2\beta} \Bigr),

$$

а коэффициенты

$A$, $B$, $C$ определяются через $u(x,y)$ точно так же, как $v(\xi,\eta)$.



Существование функции Грина"--~Адамара $G(\xi,\eta;\xi_0,\eta_0)$ 

второй задачи Дарбу для уравнения \eqref{Attfor2015_1} при $\lambda=0$

доказано Геллерстедтом \cite{Attfor2015_lit_02}:

%Она имеет вид

\begin{displaymath}

G(\xi,\eta;\xi_0,\eta_0) = 

\left\{ \begin{array}{ll}

     g (\xi,\eta;\xi_0,\eta_0), \quad \eta\geqslant\xi_0,\\

\bar{g}(\xi,\eta;\xi_0,\eta_0), \quad \eta\leqslant\xi_0.

\end{array} \right.

\end{displaymath}

Здесь $g$ "--- функция Римана, а $\bar{g}$ "--- функция, обладающая 

следующими свойствами:

\begin{itemize}

\item[1)] $\bar{g}(\xi,\eta;\xi_0,\eta_0)$ по переменным $(\xi_0,\eta_0)$ 

удовлетворяет уравнению \eqref{Attfor2015_29}, а по переменным $(\xi,\eta)$ "--- ему сопряженному\hypertarget{at:usl2}{;} 



\item[2)] $\bar{g}(\xi,\xi_0;\xi_0,\eta_0)-g(\xi,\xi_0;\xi_0,\eta_0)=$

$$\cos \pi\beta \lim\limits_{\varepsilon\rightarrow 0}

\bigl(g(\xi,\xi_0+\varepsilon;\xi_0,\xi_0+\varepsilon)g(\xi,\xi_0+\varepsilon;\xi_0,\eta_0)\bigr);$$



\item[3)] $\bar{g}$ при $\eta=\xi$ обращается в нуль порядка $2\beta$.



\end{itemize}



Так как функция Римана $g(\xi,\eta;\xi_0,\eta_0)$ удовлетворяет соотношениям

$$

g(\xi_0,\eta;\xi_0,\eta_0)=\exp \int _{\eta_0}^\eta A(\xi_0,t)dt=

\Bigl(\frac{\eta-\xi_0}{\eta_0-\xi_0}\Bigr)^\beta \exp

 \int _{\eta_0}^\eta A_1(\xi_0,t)dt,

$$

$$

g(\xi,\eta_0;\xi_0,\eta_0)=\exp \int _{\xi_0}^\xi B(t,\eta_0)dt=

\Bigl(\frac{\eta_0-\xi}{\eta_0-\xi_0}\Bigr)^\beta \exp

 \int _{\xi_0}^\xi B_1(t,\eta_0)dt,

$$

условие \hyperlink{at:usl2}{2}) можно записать в виде

\begin{multline}

\bar{g}(\xi,\xi_0;\xi_0,\eta_0)-g(\xi,\xi_0;\xi_0,\eta_0)=

\\

=\cos \pi\beta A^*(\xi,\xi_0,\eta_0)\Bigl(\frac{\xi_0-\xi}{\eta_0-\xi_0}\Bigr)

\exp\biggl(

 \int _{\xi_0}^{\eta_0} \frac{\widetilde{b}(\xi_0,t)dt}{(t-\xi_0)^{2\beta}} -

 \int _{\xi}^{\xi_0}    \frac{\widetilde{b}(t,\xi_0)dt}{(\xi_0-t)^{2\beta}}

\biggr),

\end{multline}

где

$$

A^*(\xi;\xi_0,\eta_0)=\lim\limits_{x\rightarrow 0} \exp\biggl(

 \int _{\xi}^{\xi_0} \frac{\widetilde{a}(t,\xi+\varepsilon)dt}{(\xi_0+\varepsilon-t)^{4\beta}}+

 \int _{\xi_0+\varepsilon}^{\eta_0} \frac{\widetilde{a}(\xi_0,t)dt}{(t-\xi_0)^{4\beta}}

\biggr),

$$

$$

\widetilde{a}(\xi,\eta)=\frac{1}{4} \Bigl(\frac{4}{m+2}\Bigr)^{4\beta} a

\Bigl(\frac{\xi+\eta}{2},-\Bigl(\frac{m+2}{4}\Bigr)^{1-2\beta}

(\eta-\xi)^{1-2\beta} \Bigr),

$$

$$

\widetilde{b}(\xi,\eta)=\frac{1}{4} \Bigl(\frac{4}{m+2}\Bigr)^{4\beta} b

\Bigl(\frac{\xi+\eta}{2},-\Bigl(\frac{m+2}{4}\Bigr)^{1-2\beta}

(\eta-\xi)^{1-2\beta} \Bigr).

$$



Таким образом, для существования функции Грина"--~Адамара 

необходимо потребовать, чтобы выполнялось неравенство

$$

A^*(\xi;\xi_0,\eta_0)<\infty \quad (0<\xi<1,\, 0<\xi_0,\, \eta_0<1).

$$

Это условие, впервые выписанное А.~М.~Нахушевым  \cite{Attfor2015_lit_03}, выполняется, если, например, $m<2$ или $a(x, y)=O(1)y^\alpha$, 

где $\alpha>{m}/{2}-1$ в $\bar{\Omega}$. 



\smallskip



Не нарушая общности, {\it д\,о\,к\,а\,з\,а\,т\,е\,л\,ь\,с\,т\,в\,о}\/ теоремы~\hyperlink{att:th1}{1}  проведем в случае

однородной задачи Гурса ($\varphi(y)\equiv 0$) для  неоднородного уравнения.



Пусть существует решение $u(x,y)$ однородной задачи Гурса и, как и ранее,

$$u(x , 0)=\tau(x),\quad u_y(x , 0)=\nu(x).$$ 

Тогда в силу тождества Римана для уравнения \eqref{Attfor2015_29} и свойств функции Грина второй задачи 

Дарбу существует такой интегральный оператор $A_\eta^\xi$, определенный

на функциях $\nu(x)$, имеющих дробный интеграл порядка $1-2\beta$ 

с~началом в точке 0 и с концом в любой точке $x\in{}]0 , 1]$, что для любой

точки $(x,y)\in\Omega^-$

\begin{equation}

\label{Attfor2015_30}

u(x,y)=A_\eta^\xi \nu+\mu \int _0^\xi d\xi_1  \int _{\xi_1}^\eta

\frac{A_i(\xi_1,\eta_1) D_{0x}^{\alpha_i}\tau}{(\eta_1-\xi_1)^{4\beta}}

G(\xi_1,\eta_1;\xi_1,\eta)d\eta_1+F(x,y),

\end{equation}

где

$$

F(x,y)= \int _0^\xi d\xi_1  \int _{\xi_1}^\eta

\frac{f(\xi_1,\eta_1)}{(\eta_1-\xi_1)^{4\beta}}

G(\xi_1,\eta_1;\xi,\eta)d\eta_1.

$$



Из \ref{Attfor2015_30} при  $\xi=\eta=x$ и $y\rightarrow 0$ имеем

\begin{equation}

\label{Attfor2015_31}

\tau(x)=A_x^x \nu+\mu \int _0^x d\xi_1  \int _{\xi_1}^x

\frac{A(\xi_1,\eta_1) D_{0x}^{\alpha_i}\tau}{(\eta_1-\xi_1)^{4\beta}}

\bar{g}(\xi_1,\eta_1;x,x)d\eta_1+F(x,x).

\end{equation}



Допустим теперь, что $(x,y)\in\Omega^+$. 

В~уравнении \eqref{Attfor2015_1} произведем преобразования $\widetilde{x}=x$, $\widetilde{y}=-y$.

В~результате область $\Omega^+$ отобразится в область $\widetilde{\Omega}^-$,

ограниченную отрезком 

$$\bar{J}=\{(\widetilde{x},\widetilde{y}): 0\leqslant\widetilde{x}\leqslant 1, \widetilde{y}=0\}$$ \\

прямой $\widetilde{y}=0$ и характеристиками

$$

\widetilde{x}-\frac{2}{m+2} (-\widetilde{y})^{\frac{m+2}{2}}=0, \quad

\widetilde{x}+\frac{2}{m+2} (-\widetilde{y})^{\frac{m+2}{2}}=1

$$

уравнения

\begin{multline}

\label{Attfor2015_32}

\widetilde{u}_{\widetilde{y}\widetilde{y}}-

(-\widetilde{y})^m \widetilde{u}_{\widetilde{x}\widetilde{x}}+

     a (\widetilde{x},-\widetilde{y})\widetilde{u}_{\widetilde{x}}-

\bar{b}(\widetilde{x},-\widetilde{y})\widetilde{u}_{\widetilde{y}}+

     c (\widetilde{x},-\widetilde{y})\widetilde{u}=

\\

=\lambda a_i (\widetilde{x},-\widetilde{y}) 

D_{0\widetilde{\xi}}^{\alpha_i}\, u(\widetilde{\xi} , 0)+

d(\widetilde{x},-\widetilde{y}).

\end{multline}



\noindent

Здесь

$$

\widetilde{u}=\widetilde{u}(\widetilde{x},\widetilde{y})=u(x,-y), \quad 

\widetilde{\xi}=\widetilde{x}-\frac{2}{m+2} (-\widetilde{y})^{\frac{m+2}{2}}.

$$



Так как для всех $(\widetilde{x},\widetilde{y})\in\widetilde{\Omega}^-$

$$a(\widetilde{x},-\widetilde{y})= a(\widetilde{x},\widetilde{y}),\quad  

b(\widetilde{x},-\widetilde{y})=-b(\widetilde{x},\widetilde{y}),\quad  

c(\widetilde{x},-\widetilde{y})= c(\widetilde{x},\widetilde{y}),$$ 

уравнение \eqref{Attfor2015_32} можно переписать в форме

\begin{multline}

\label{Attfor2015_33}

\widetilde{u}_{\widetilde{y}\widetilde{y}}-

(-\widetilde{y})^m \widetilde{u}_{\widetilde{x}\widetilde{x}}+

a(\widetilde{x},\widetilde{y})\widetilde{u}_{\widetilde{x}}+

b(\widetilde{x},\widetilde{y})\widetilde{u}_{\widetilde{y}}+

c(\widetilde{x},\widetilde{y})\widetilde{u}=

\\

=\lambda a_i (\widetilde{x},\widetilde{y}) 

D_{0\widetilde{\xi}}^{\alpha_i}\, \widetilde{u}(\widetilde{\xi} , 0)+

\widetilde{d}(\widetilde{x},\widetilde{y}).

\end{multline}



Очевидно, что

\begin{equation}

\label{Attfor2015_34}

\tau(\widetilde{x})=\tau(x),\quad \lim\limits_{\widetilde{y}\rightarrow 0^-}

\widetilde{u}_{\widetilde{y}}(\widetilde{x},\widetilde{y})=-\nu(\widetilde{x}),

\end{equation}

$$

\widetilde{u} \Bigl(

\frac{2}{m+2} (-\widetilde{y})^{\frac{m+2}{2}},\widetilde{y} 

\Bigr)=0,

\qquad

-\Bigl(\frac{m+2}{2}\Bigr)^{\frac{2}{m+2}}\leqslant\widetilde{y}\leqslant 0.

$$

Пусть, как и выше,

$$

\widetilde{f}(\widetilde{\xi},\widetilde{\eta})=

-\frac{1}{4} 

\Bigl(\frac{4}{m+2}\Bigr)^{4\beta}\widetilde{d}

\Bigl(

\frac{\widetilde{\xi}+\widetilde{\eta}}{2},-\Bigl(\frac{m+2}{4}\Bigr)^{1-2\beta}

(\widetilde{\eta}-\widetilde{\xi})^{1-2\beta}

\Bigr),

$$

$$

\widetilde{F}(\widetilde{x},\widetilde{y})= \int _0^{\widetilde{\xi}}

 d\xi_1  \int _{\xi_1}^{\widetilde{\eta}}

\frac{\widetilde{f}(\xi_1,\eta_1)}{(\eta_1-\xi_1)^{4\beta}}

G(\xi_1,\eta_1;\xi,\eta)d\eta_1,

$$

$$

\widetilde{A_i}(\widetilde{\xi},\widetilde{\eta})=

\frac{1}{4} 

\Bigl(\frac{4}{m+2}\Bigr)^{4\beta}\widetilde{a_i}

\Bigl(

\frac{\widetilde{\xi}+\widetilde{\eta}}{2},

-\Bigl(\frac{m+2}{4}\Bigr)^{1-2\beta}

(\widetilde{\eta}-\widetilde{\xi})^{1-2\beta}

\Bigr),

$$

где 

$$

\widetilde{\xi} =\widetilde{x}-\frac{2}{m+2} (-\widetilde{y})^{\frac{m+2}{2}}, \quad

\widetilde{\eta}=\widetilde{x}+\frac{2}{m+2} (-\widetilde{y})^{\frac{m+2}{2}}.

$$

Тогда в силу \eqref{Attfor2015_30} имеем

\begin{equation}

\label{Attfor2015_35}

\widetilde{u}(\widetilde{x},\widetilde{y})=-A_{\widetilde{\eta}}^{\widetilde{\xi}}\,\nu+

\mu \int _0^{\widetilde{\xi}} d\xi_1  \int _{\xi_1}^{\widetilde{\eta}}

\frac{\widetilde{A}_i(\xi_1,\eta_1) D_{0\xi_1}^{\alpha_i}\tau}{(\eta_1-\xi_1)^{4\beta}}

G(\xi_1,\eta_1;\widetilde{\xi},\widetilde{\eta})d\eta_1+

\widetilde{F}(\widetilde{x},\widetilde{y}).

\end{equation}

Равенство \eqref{Attfor2015_35} дает нам право записать

\begin{multline}

\label{Attfor2015_36}

\tau(x)=-A_x^x \nu+\mu \int _0^x d\xi_1  \int _{\xi_1}^x

\frac{\widetilde{A}_i(\xi,\eta_1) D_{0\xi_1}^{\alpha_i}\tau}{(\eta_1-\xi_1)^{4\beta}}

\bar{g}(\xi_1,\eta_1;x,x)d\eta_1 + \\ +\widetilde{F}(x,x), \quad  0\leqslant x\leqslant 1.

\end{multline}



Равенство \eqref{Attfor2015_36} выражает фундаментальное соотношение между

$\tau(x)$ и~$\nu(x)$, принесенное из области $\Omega^+$ на линию 

параболического выражения~${y=0}$.



Из \eqref{Attfor2015_31} и \eqref{Attfor2015_36} после их почленного  сложения находим

\begin{equation}

\label{Attfor2015_37}

\tau(x)=\mu \int _0^x d\xi_1  \int _{\xi_1}^x

\frac{\bar{A}_i(\xi_1,\eta_1) D_{0\xi_1}^{\alpha_i}\tau}{(\eta_1-\xi_1)^{4\beta}}

\bar{g}(\xi_1,\eta_1;x,x)d\eta_1+\bar{F}(x),

\end{equation}

где

$$

\bar{A}_i(\xi,\eta)=\frac{1}{2}\bigl(A_i(\xi,\eta)+\widetilde{A}_i(\xi,\eta)\bigr),

\quad 

\bar{F}(x)=\frac{1}{2}\bigl( F(x,x)+\widetilde{F}(x,x)\bigr).

$$



Рассмотрим оператор

\begin{equation}

\label{Attfor2015_38}

V_i \tau= \int _0^x d\xi_1  \int _{\xi_1}^x

\frac{A_i(\xi_1,\eta_1) D_{0\xi_1}^{\alpha_i}\tau}{(\eta_1-\xi_1)^{4\beta}}

\bar{g}(\xi_1,\eta_1;x,x)d\xi_1.

\end{equation}

Очевидно, что 

$$V_\tau=\sum\limits_{i=1}^n V_i \tau,\qquad  

V_i \tau= \int _0^x  D_{0\xi_1}^{\alpha_i}\tau\, G_i(x,\xi_1)d\xi_1,

$$

где

$$

G_i(x,\xi_1)= \int _{\xi_1}^x 

\frac{\bar{A}_i(\xi_1,\eta_1) \bar{g}(\xi_1,\eta_1;x,x)}{(\eta_1-\xi_1)^{4\beta}}d\eta_1.

$$



Не нарушая общности, можно предположить, 

что $\alpha_i\leqslant0$ при $i=1 , 2,\ldots,m$ и $\alpha_i>0$ 

при $i=m+1, m+2, \ldots,n$.



Пусть $i\geqslant m+1$. Тогда для любой функции $\tau(\xi)\in D(V_i)$ на 

основании определения производной дробного порядка

\begin{multline}

V_i \tau= \int _0^x \frac{G_i(x,t)}{\Gamma(1-\alpha_i)}

\biggl(\dfrac{d}{dt} 

 \int _0^t \frac{\tau(\eta)}{(t-\eta)^{\alpha_i}} d\eta\biggr) dt=

\\

=\biggl( \frac{G_i(x,t)}{\Gamma(1-\alpha_i)} \int _0^t

\frac{\tau(\eta)d\eta}{(t-\eta)^{\alpha_i}}

\biggr)\bigg|_{t=0}^{t=x} -

 \int _0^x \frac{\frac{\partial}{\partial t}G_i(x,t)}{\Gamma(1-\alpha_i)}dt 

 \int _0^t \frac{\tau(\eta)}{(t-\eta)^{\alpha_i}} d\eta.

\end{multline}



Отсюда после законной перестановки порядка интегрирования в двойном интеграле получаем

\begin{multline}

V_i \tau= \frac{G_i(x,x)}{\Gamma(1-\alpha_i)} \int _0^x 

\frac{\tau(\eta)}{(x-\eta)^{\alpha_i}} d\eta-

\\

- \int _0^x \tau(\eta)d\eta  \int _t^x 

\frac{\frac{\partial}{\partial t}G_i(x,t)}{\Gamma(1-\alpha_i)(t-\eta)^{\alpha_i}}dt, 

\quad \forall i\geqslant m+1,

\end{multline}

или

$$

V_i \tau=G_i(x,x)D_{0x}^{\alpha_i-1}\tau(\eta)+

 \int _0^x G_i^0(x,\eta)\tau(\eta)d\eta,

$$

где

$$

G_i^0(x,\eta)= \int _t^x 

\frac{\frac{\partial}{\partial t}G_i(x,t)}{\Gamma(1-\alpha_i)(t-\eta)^{\alpha_i}}dt.

$$

На основании \eqref{Attfor2015_4} $G_i^0(x,\eta)\in C(\bar{\Omega})$.



Пусть теперь $i\leqslant m$, тогда, согласно определению дробного интеграла порядка 

$-\alpha_i$, имеем

$$

V_i \tau= \int _0^x \frac{G_i(x,t)}{\Gamma(-\alpha_i)} dt \int _0^t 

\frac{\tau(\eta)}{(t-\eta)^{\alpha_i+1}} d\eta.

$$



Формула перестановки Дирихле говорит о том, что это равенство эквивалентно равенству

$$

V_i \tau= \int _0^x G_i^0(x,\eta)\tau(\eta)d\eta,

\quad i=1 , 2,\ldots,m,

$$

где

$$

G_i^0(x,\eta)= \int _\eta^x 

\frac{G_i(x,t)}{\Gamma(-\alpha_i)(t-\eta)^{\alpha_i+1}}dt.

$$

Очевидно, что $G_i^0(x,\eta)\in C(\bar{\Omega})$.



Таким образом, доказано, что $V$ "--- оператор, действующий на $\tau\in C[0 , 1]$ по формуле 

\eqref{Attfor2015_38}, представим в виде

$$

V \tau= \int _0^x K^*(x,t)\tau(t)dt,

$$

где ядро $K^*(x,t)$ при $x=t$ может иметь лишь слабую особенность.



Следовательно, уравнение \eqref{Attfor2015_37} эквивалентно интегральному  уравнению Вольтерра второго рода

\begin{equation}

\label{Attfor2015_39}

\tau(x)=\mu \int _0^x K^*(x,t)\tau(t)dt+\bar{F}(x).

\end{equation}



После того как из \eqref{Attfor2015_39} $\tau(x)$ найдено однозначно,  решение $u(x,y)$ задачи Гурса \eqref{Attfor2015_3} для уравнения  \eqref{Attfor2015_1} определяется как решение этой же задачи, но уже для ненагруженного уравнения

$$

u_{yy}-|y|^m u_{xx}+a u_x+b u_y+c u=d_0(x,y),

$$ 

где

$$

d_0(x,y)=\lambda a_i D_{0\xi}^{\alpha_i} \tau(\xi)+d(x,y).

$$





\thanks{{\bf ORCID}

{



{\bf Анатолий Хусеевич Аттаев:} \url{http://orcid.org/0000-0001-5864-6283}

}

}



\pagebreak



\RusRef



\begin{thebibliography}{99}



\RBibitem{litAtt01}

\by Нахушев~А.~М.

\book Нагруженные уравнения и их применение

\yr 2012

\publ Наука

\publaddr М.

\totalpages 232

\elib{http://elibrary.ru/item.asp?id=20886619}





\RBibitem{litAtt02}

\by Нахушев~А.~М.

\book Уравнения математической биологии

\yr 1995

\publ Высшая школа

\publaddr М.

\totalpages 301

\zmath{https://zbmath.org/?q=an:01785571}





\RBibitem{litAtt03}

\by Нахушев~А.~М.

\paper О задаче Дарбу для одного вырождающегося нагруженного интегро-дифференциального уравнения второго порядка

\jour Дифференц. уравнения

\yr 1976

\vol 12

\issue 1

\pages 103--108

\mathnet{http://mi.mathnet.ru/rus/de2654}

%\mathscinet{http://www.ams.org/mathscinet-getitem?mr=1670100}

\zmath{https://zbmath.org/?q=an:0349.45011}



\RBibitem{litAtt04}

\by Казиев~М.~В.

\paper Задача Гурса для одного нагруженного интегро-дифференциального уравнения

\jour Дифференц. уравнения

\yr 1981

\vol 17

\issue 2

\pages 313--319

\mathscinet{http://www.ams.org/mathscinet-getitem?mr=606821}

\zmath{https://zbmath.org/?q=an:0487.45009}



\RBibitem{litAtt05}

\by Репин~О.~А., Тарасенко~А.~В.

\paper Задача Гурса и Дарбу для одного нагруженного интегро-дифференциального уравнения второго порядка

\jour Математический журнал. Алматы

\yr 2011

\vol 11

\issue 2

\pages 64--72

\elink \url{http://www.math.kz/images/journal/2011-2/Repin_Tarasenko.pdf} (дата обращения: 23.10.2015)

\mathscinet{}



\RBibitem{litAtt06}

\by Аттаев~А.~Х.

\paper Задача Гурса для локально нагруженного уравнения со степенным параболическим вырождением

\jour Доклады Адыгской (Черкесской) Международной академии наук

\yr 2008

\vol 10

\issue 2

\pages 14--17

\elib{http://elibrary.ru/item.asp?id=12792278}



\RBibitem{litAtt07}

\by Аттаев~А.~Х.

\paper Задача Гурса для нагруженного гиперболического уравнения

\jour Доклады Адыгской (Черкесской) Международной академии наук

\yr 2014

\vol 16

\issue 3

\pages 9--12

\elib{http://elibrary.ru/item.asp?id=22543031}



\RBibitem{litAtt08}

\by Огородников~Е.~Н.

\paper Некоторые характеристические задачи для систем нагруженных дифференциальных уравнений и их связь с~нелокальными краевыми задачами

\jour Вестн. Сам. гос. техн. ун-та. Сер. Физ.-мат. науки

\yr 2003

\issue 19

\pages 22--28

\mathnet{http://mi.mathnet.ru/rus/vsgtu134}

\crossref{10.14498/vsgtu134}



\RBibitem{litAtt09}

\by Огородников~Е.~Н.

\paper Корректность задачи Коши--Гурса для системы вырождающихся нагруженных гиперболических уравнений в~некоторых специальных случаях и ее равносильность задачам с~нелокальными краевыми условиями

\jour Вестн. Сам. гос. техн. ун-та. Сер. Физ.-мат. науки

\yr 2004

\issue 26

\pages 26--38

\mathnet{http://mi.mathnet.ru/rus/vsgtu173}

\crossref{10.14498/vsgtu173}



\RBibitem{Attfor2015_lit_01}

\by Нахушев~А.~М.

\paper О нелокальных краевых задачах со смещением и их связи с нагруженными уравнениями

\jour Дифференц. уравнения

\yr 1985

\vol 21

\issue 1

\pages 92--101

\mathscinet{http://www.ams.org/mathscinet-getitem?mr=777785}

\zmath{https://zbmath.org/?q=an:0573.35066}



\Bibitem{Attfor2015_lit_02}

\by Gellerstedt~S.

\paper Sur une \'equation lin\'eaire aux d\'eriv\'ees partielles de type mixte

\jour Ark. Mat. Astron. Fys. A 

\vol 25

\yr 1937

\issue 29

\pages 1--23

\zmath{https://zbmath.org/?q=an:0016.21203}



\RBibitem{Attfor2015_lit_03}

\by Нахушев~А.~М.

\book Об одном классе линейных краевых задач для гиперболического и смешанного типов уравнений второго порядка

\yr 1992

\publ Эльбрус

\publaddr Нальчик

\totalpages 154





\end{thebibliography}





\RuDate



\newpage



\makeentitle



\thanks{{\bf ORCID}

{

{\bf Anatoly H. Attaev:} \url{http://orcid.org/0000-0001-5864-6283}

}

}







\EngRef



\begin{thebibliography}{99}



\Bibitem{litAtt01:eng}

\lang In Russian

\by Nakhushev~A.~M.

\book Nagruzhennye uravneniia i ikh primenenie \rm [Loaded equations and its applications]

\yr 2012

\publ Nauka

\publaddr Moscow

\totalpages 232





\Bibitem{litAtt02:eng }

\lang In Russian

\by Nakhushev~A.~M.

\book Uravneniia matematicheskoi biologii \rm [Equations of mathematical biology]

\yr 1995

\publ Vysshaia shkola

\publaddr Moscow

\totalpages 301





\Bibitem{litAtt03:eng}

\lang In Russian

\by Nakhushev~A.~M.

\paper On the Darboux problem for a nondegenerate loaded integro-differential equation of the second order

\jour Differ. Uravn.

\yr 1976

\vol 12

\issue 1

\pages 103--108





\Bibitem{litAtt04:eng}

\by Kaziev~V.~M.

\paper Goursat's problem for a loaded integrodifferential equation

\jour Differ. Equations 

\vol 17

\issue 2

\pages 216--220 

\yr 1981

\mathscinet{http://www.ams.org/mathscinet-getitem?mr=606821}

\zmath{https://zbmath.org/?q=an:0487.45009}



\Bibitem{litAtt05:eng}

\lang In Russian

\by Repin~O.~A., Tarasenko~A.~V.

\paper The Goursat and Darboux problems investigated for a loaded integro-differential equation of the second order

\jour Mathematical Journal

\yr 2011

\vol 11

\issue 2

\pages 64--72

\elink Retrieved from 

 \url{http://www.math.kz/images/journal/2011-2/Repin_Tarasenko.pdf} (October  10, 2015) 





\Bibitem{litAtt06:eng}

\lang In Russian

\by Attaev~A.~H.

\paper A Goursat problem for a locally loaded parabolic equations with power degeneration

\jour Doklady Adygskoi (Cherkesskoi) Mezhdunarodnoi akademii nauk

\yr 2008

\vol 10

\issue 2

\pages 14--17





\Bibitem{litAtt07:eng}

\lang In Russian

\by Attaev~A.~H.

\paper A Goursat problem for a loaded hyperbolic equation

\jour Doklady Adygskoi (Cherkesskoi) Mezhdunarodnoi akademii nauk

\yr 2014

\vol 16

\issue 3

\pages 9--12





\Bibitem{litAtt08:eng}

\lang In Russian

\by Ogorodnikov~E.~N.

\paper Some characteristic problems for loaded systems of differential equations and their relationship with non-local boundary value problems

\jour Vestn. Samar. Gos. Tekhn. Univ. Ser. Fiz.-Mat. Nauki \rm [J. Samara State Tech. Univ., Ser. Phys. \& Math. Sci.]

\yr 2003

\issue 19

\pages 22--28

\crossref{10.14498/vsgtu134}



\Bibitem{litAtt09:eng}

\lang In Russian

\by Ogorodnikov~E.~N.

\paper The correctness of the Cauchy--Goursat problem for loaded degenerate hyperbolic equations  in some special cases, and its equivalent to the problem with nonlocal boundary conditions

\jour Vestn. Samar. Gos. Tekhn. Univ. Ser. Fiz.-Mat. Nauki \rm [J. Samara State Tech. Univ., Ser. Phys. \& Math. Sci.]

\yr 2004

\issue 26

\pages 26--38

\crossref{10.14498/vsgtu173}



\Bibitem{Attfor2015_lit_01:eng}

\by Nakhushev~A.~M.

\paper Nonlocal boundary problems with displacement and their relation to loaded equations

\jour Differ. Equations

\vol 21

\issue 1

\pages 74--81 

\yr 1985

\mathscinet{http://www.ams.org/mathscinet-getitem?mr=777785}

\zmath{https://zbmath.org/?q=an:0573.35066}



\Bibitem{Attfor2015_lit_02:eng}

\by Gellerstedt~S.

\paper Sur une \'equation lin\'eaire aux d\'eriv\'ees partielles de type mixte

\jour Ark. Mat. Astron. Fys. A 

\vol 25

\yr 1937

\issue 29

\pages 1--23

\zmath{https://zbmath.org/?q=an:0016.21203}



\Bibitem{Attfor2015_lit_03:eng}

\lang In Russian

\by Nakhushev~A.~M.

\book Ob odnom klasse lineinykh kraevykh zadach dlia giperbolicheskogo i smeshannogo tipov uravnenii vtorogo poriadka \rm [One Class of Linear Boundary-Value Problems for Hyperbolic and Mixed-Type Equations of the Second Order]

\yr 1992

\publ El'brus

\publaddr Nal'chik

\totalpages 154







\end{thebibliography}



\EngDate





\newpage

\Russian





%\end{document}

